% Source: ME14_v1.html
\documentclass[11pt]{article}
\usepackage[margin=1in]{geometry}
\usepackage{amsmath}
\usepackage{amssymb}
\usepackage{siunitx}
\usepackage{enumitem}

\title{ME14: Oscillations}
\author{}
\date{}

\begin{document}
\maketitle

\section*{ME14: Oscillations}

\begin{enumerate}[leftmargin=*,label=\textbf{Question \arabic*.}]
\item \hfill \textit{/ 10 pts}

Suppose a horizontal block-spring system has a spring constant 1,055 N/m and block of mass 11 kg. Calculate the angular frequency in rad/s of its harmonic motion..

\textit{Correct answer: } 9.7933

\item \hfill \textit{/ 10 pts}

Suppose a horizontal block-spring system has a spring constant of 1240 N/m and a block of mass 15.0 kg. Calculate:

\begin{itemize}

\item the angular frequency in rad/s: [Select]

\item the frequency in Hz: [Select]
, and

\item the period in s: [Select]
.

\end{itemize}

\textit{Answer 1:}
\begin{itemize}
  \item 13.5 read/s
  \item 8.42 rad/s
  \item \textbf{[correct]} 9.09 rad/s
  \item 11.5 rad/s
\end{itemize}
\textit{Answer 2:}
\begin{itemize}
  \item 1.34 Hz
  \item 1.83 Hz
  \item 2.15 Hz
  \item \textbf{[correct]} 1.45 Hz
\end{itemize}
\textit{Answer 3:}
\begin{itemize}
  \item 0.465 s
  \item \textbf{[correct]} 0.691 s
  \item 0.546 s
  \item 0.676 s
\end{itemize}

\item \hfill \textit{/ 10 pts}

A tunnel is drilled straight through the center of an airless planet to the other side. For simplicity, the planet is assumed to have constant density. The equation for the motion of an object falling through the tunnel is given by

$\frac{d^2r}{dt^2}+\left(\frac{4\pi G\rho}{3}\right)r=0$


where \emph{G} = 6.67 × 10$^{-11}$ m$^{3}$/(kgᐧ݀s$^{2}$) is the gravitation constant. The planet's density and radius are 𝜌 = 4,740 kg/m$^{3}$ and \emph{R} = 1,775 km.

If an astronaut jumps in the tunnel, how many \textbf{\emph{hours}} would it take for her to fall all the way to the other side and return to her original position?



\emph{T} = \rule{2cm}{0.4pt} h (Note the units!)

\textit{Correct answer: } 1.5166

\item \hfill \textit{/ 10 pts}

Suppose a harmonic oscillator is described by \emph{ma$_{z}$ + kz} = 0 where \emph{m} = 11.7 kg and \emph{k} = 94 N/m. Calculate the period of the harmonic motion.

\textit{Correct answer: } 2.2167

\item \hfill \textit{/ 10 pts}

Fill In the Blanks with numerical values. \emph{\textbf{Give three digits for each answer}}, rounding as appropriate. For example, if your answer is "2", report "2.00". If your answer is "1.5555", report "1.56". If your answer is "0.7", report "0.700".



Suppose a mass-spring system has the description \emph{x}(\emph{t}) = 4.00 cos(2.00\emph{t} - 0.500), where each number has unshown MKS units.



The amplitude is
 m, the angular frequency is
 rad/s, the phase shift is
 rad, the frequency is
 Hz; the period is
 s.

\textit{Answer 1:}
\begin{itemize}
  \item \textbf{[correct]} 4.00
\end{itemize}
\textit{Answer 2:}
\begin{itemize}
  \item \textbf{[correct]} 2.00
\end{itemize}
\textit{Answer 3:}
\begin{itemize}
  \item \textbf{[correct]} -0.500
  \item \textbf{[correct]} -.500
\end{itemize}
\textit{Answer 4:}
\begin{itemize}
  \item \textbf{[correct]} 0.318
  \item \textbf{[correct]} .318
\end{itemize}
\textit{Answer 5:}
\begin{itemize}
  \item \textbf{[correct]} 3.14
\end{itemize}

\item \hfill \textit{/ 10 pts}

Choose the correct answers.



Suppose a simple harmonic oscillator solution for a mass-spring system is given by \emph{x}(\emph{t}) = 0.7 cos(8\emph{t} + 1.9), where all MKS units are assumed.



For this solution,

\textbf{(a) }the amplitude is [Select]
 m;

\textbf{(b)} the angular frequency is [Select]
 rad/s;

\textbf{(c) }the phase shift is [Select]
 rad;

\textbf{(d)} the frequency is [Select]
 Hz;

\textbf{(e)} the period is [Select]
 s;

\textbf{(f)} the maximum speed is [Select]
 m/s;

\textbf{(g)} the maximum acceleration is [Select]
 m/s$^{2}$;

\textbf{(h)} the position at t = 0 s is [Select]
 m;

\textbf{(i)} the velocity at t = 0 s is [Select]
 m/s;

\textbf{(j)} the acceleration at t = 0 s is [Select]
 m/s$^{2}$;

\textbf{(k)} Find the first time t > 0 that the mass reaches x = 0. [Select]

\textbf{(l)} Find the first time t > 0 that the speed is 0. 0.155

\textit{Answer 1:}
\begin{itemize}
  \item 5.6
  \item 4.21
  \item 8
  \item .49
  \item \textbf{[correct]} 0.7
\end{itemize}
\textit{Answer 2:}
\begin{itemize}
  \item 0.7
  \item 1.27
  \item 1.9
  \item \textbf{[correct]} 8
\end{itemize}
\textit{Answer 3:}
\begin{itemize}
  \item \textbf{[correct]} 1.9
  \item 0.7
  \item 8
\end{itemize}
\textit{Answer 4:}
\begin{itemize}
  \item \textbf{[correct]} 1.27
  \item 1.58
  \item 0.785
  \item 8
  \item 1.9
\end{itemize}
\textit{Answer 5:}
\begin{itemize}
  \item \textbf{[correct]} 0.785
  \item 0.672
  \item 1.86
  \item 1.27
\end{itemize}
\textit{Answer 6:}
\begin{itemize}
  \item 2.8
  \item 4.2
  \item 3.6
  \item \textbf{[correct]} 5.6
\end{itemize}
\textit{Answer 7:}
\begin{itemize}
  \item 32.5
  \item 22.4
  \item 18.6
  \item \textbf{[correct]} 44.8
\end{itemize}
\textit{Answer 8:}
\begin{itemize}
  \item 0.6996
  \item -0.523
  \item \textbf{[correct]} -0.226
  \item -0.334
\end{itemize}
\textit{Answer 9:}
\begin{itemize}
  \item -4.15
  \item \textbf{[correct]} -5.30
  \item -2.62
  \item 2.62
  \item 5.30
\end{itemize}
\textit{Answer 10:}
\begin{itemize}
  \item -13.2
  \item 19.3
  \item 8.92
  \item \textbf{[correct]} 14.48
  \item -5.80
\end{itemize}
\textit{Answer 11:}
\begin{itemize}
  \item \textbf{[correct]} 0.352
  \item 0.265
  \item 0.155
  \item 0.785
  \item 0.418
\end{itemize}
\textit{Answer 12:}
\begin{itemize}
  \item 0.405
  \item \textbf{[correct]} 0.155
  \item 0.261
  \item 0.352
\end{itemize}

\item \hfill \textit{/ 10 pts}

A mass spring system has an initial speed of 2.6 m/s in the negative x-direction and an initial position of 0.51 m. Calculate the phase shift in radians, assuming a cosine solution, with a spring constant of 539 N/m and a mass of 4 kg.

\textit{Correct answer: } 0.4138

\item \hfill \textit{/ 10 pts}

A simple harmonic oscillator consists of a block of an unknown mass in kilograms and spring with spring constant 514 N/m on a frictionless horizontal surface. The block is moved a distance 1.6 mfrom equilibrium and released from rest. If the block's maximum speed at the equilibrium point is 12.9 m/s, calculate the mass of the block in kg.

\textit{Correct answer: } 7.9072

\item \hfill \textit{/ 10 pts}

A simple harmonic oscillator consists of a block of mass 4.7 kg and spring with spring constant 183 N/m on a frictionless horizontal surface. The block is moved from equilibrium and released from rest. If the block's maximum speed 10 m/s, calculate the amplitude \emph{A} of the motion in meters.



\emph{A} = \rule{2cm}{0.4pt} m

\textit{Correct answer: } 1.6026

\item \hfill \textit{/ 10 pts}

A 7-kg block attached to a spring of force constant 133 N/m undergoes simple harmonic motion. At a particular time the mass' speed is 2.9 m/s and its position is 0.26 m. Calculate the amplitude of its oscillation.



\emph{A }= \rule{2cm}{0.4pt} m

\textit{Correct answer: } 0.7143

\item \hfill \textit{/ 10 pts}

A 1.5-kg block attached to a spring oscillates with an energy of 5.8 J. Calculate the block's speed when its displacement is 14\% of the oscillation amplitude.



\emph{v} = \rule{2cm}{0.4pt} m/s

\textit{Correct answer: } 2.7535

\item \hfill \textit{/ 10 pts}

Calculate the length of a simple pendulum on Earth if its angular frequency is 2.82 rad/s.



 \emph{L} = \rule{2cm}{0.4pt} m

\textit{Correct answer: } 1.2323

\item \hfill \textit{/ 10 pts}

Calculate the length in cm of a simple pendulum on Earth if its period is 1.579 s.



\emph{L} = \rule{2cm}{0.4pt} cm

\textit{Correct answer: } 61.8914

\item \hfill \textit{/ 10 pts}

An object is free to oscillate around a pivot located a distance 1.41 m from its center of mass. Calculate the period of small oscillations if its mass is 11.5 kg and its moment of inertia about the pivot is 15 kgᐧm$^{2}$.



\emph{T} = \rule{2cm}{0.4pt} s

\textit{Correct answer: } 1.9304

\item \hfill \textit{/ 10 pts}

An object is free to oscillate around a pivot that is a distance 0.48 m from its center of mass. Calculate the frequency of small oscillations if its mass is 6 kg and its moment of inertia is about the pivot is 8.6 kgᐧm$^{2}$.



\emph{f} = \rule{2cm}{0.4pt} Hz

\textit{Correct answer: } 0.2883

\item \hfill \textit{/ 10 pts}

A simple pendulum is constructed by attaching a string of negligible mass to a small, dense weight that can be considered a point particle. To obtain a longer period, use a [Select]
 string. A longer period [Select]
 obtained by using a lighter mass. A longer period also results if the pendulum is taken to an area with a [Select]
 acceleration of gravity.

\textit{Answer 1:}
\begin{itemize}
  \item \textbf{[correct]} longer
  \item higher tension
  \item shorter
\end{itemize}
\textit{Answer 2:}
\begin{itemize}
  \item can be
  \item \textbf{[correct]} can't be
\end{itemize}
\textit{Answer 3:}
\begin{itemize}
  \item larger
  \item \textbf{[correct]} smaller
\end{itemize}

\end{enumerate}

\end{document}
